\documentclass[a4paper,10pt]{article}
\usepackage{amsmath,amssymb}
\usepackage{fontspec}
\setmainfont{Liberation Serif}
\usepackage{graphicx}
\usepackage{eso-pic}
\usepackage{tikz}
\usepackage[most]{tcolorbox}
\usepackage{geometry}
\usepackage{fancyhdr}
\usepackage{lastpage}
\usepackage{booktabs}
\usepackage{listings}
\usepackage{enumitem}
\usepackage{microtype}

\geometry{
    a4paper,
    top=1.5cm,
    bottom=1.8cm,
    left=1.2cm,
    right=1.2cm
}

% Persistent Unremovable Canvas Watermark on Every Page
\newsavebox{\watermarkbox}
\sbox{\watermarkbox}{%
    \begin{tabular}{c}
        \includegraphics[width=8cm]{/home/johan/GATE/GATE_CS_PYQ_DB/assets/studiva-butterfly-icon.jpg} \\[-0.5em]
        {\fontsize{18}{22}\selectfont\bfseries\scshape www.studiva.co.in}
    \end{tabular}%
}
\AddToShipoutPictureBG{%
    \begin{tikzpicture}[remember picture, overlay]
        \node[opacity=0.12] at (current page.center) {\usebox{\watermarkbox}};
    \end{tikzpicture}%
}

% Header and Footer
\pagestyle{fancy}
\fancyhf{}
\fancyhead[L]{\small\scshape GATE 2018--2025 CS/IT $\cdot$ \textbf{Algorithms}}
\fancyhead[R]{\small\scshape Practice Question Bank with Solutions}
\fancyfoot[L]{\scriptsize\textsc{Copyright \copyright\ 2026} Studiva\textsuperscript{TM}}
\fancyfoot[C]{\small Page \thepage\ of \pageref{LastPage}}
\fancyfoot[R]{\scriptsize\texttt{www.studiva.co.in}}
\renewcommand{\headrulewidth}{0.4pt}
\renewcommand{\footrulewidth}{0.4pt}

% Code Listing Style (Monochrome)
\lstset{
    basicstyle=\footnotesize\ttfamily,
    numbers=none,
    frame=none,
    tabsize=2,
    breaklines=true,
    showstringspaces=false
}

\begin{document}

% ==================== PAGE 1: IMPRINT & TITLE PAGE ====================
\thispagestyle{empty}
\vspace*{0.5cm}
\begin{center}
    \includegraphics[width=2.5cm]{/home/johan/GATE/GATE_CS_PYQ_DB/assets/studiva-butterfly-icon.jpg}\\[0.6em]
    {\LARGE \textbf{GATE 2018--2025 COMPUTER SCIENCE \& IT}}\\[0.3em]
    {\Large \textbf{PREVIOUS YEAR QUESTION BANK WITH SOLUTIONS}}\\[0.4em]
    {\large \textsc{\textbf{Subject Wise: Algorithms}}}\\[0.2em]
    {\normalsize \textsc{Combined 2018--2025 Examination Papers $\cdot$ Comprehensive Edition ($\mathbf{29}$ Questions)}}\\[1.2em]
    \rule{0.6\textwidth}{0.6pt}
\end{center}

\vspace{1em}
\noindent\begin{tabular}{ll}
\textbf{Author / Compiler:} & Rupam Ghosh \\
\textbf{Publisher \& Platform:} & Studiva\textsuperscript{TM} (\texttt{www.studiva.co.in}) \\
\textbf{Official Inquiries:} & \texttt{contact@studiva.co.in} \\
\textbf{Edition:} & First Standard Edition (2026) \\
\textbf{Coverage:} & Complete GATE CS/IT Papers (2018, 2022, 2023, 2024 CS1, 2024 CS2, 2025 CS1, 2025 CS2) \\
\textbf{Syllabus Reference:} & Official GATE 2026 Syllabus (Organizing Institute: IIT Guwahati) \\
\end{tabular}

\vspace{1.5em}
\begin{tcolorbox}[enhanced, standard jigsaw, opacityback=0, colframe=black!40, arc=1mm, boxrule=0.5pt, left=10pt, right=10pt, top=8pt, bottom=8pt]
\textbf{\large COPYRIGHT \& INTELLECTUAL PROPERTY NOTICE}\\[0.6em]
Copyright \copyright\ 2026 Rupam Ghosh \& Studiva\textsuperscript{TM}. All Rights Reserved.\\[0.5em]
This educational compilation, problem organization, taxonomy classification, LaTeX typesetting, hints, and explanatory solutions are the proprietary intellectual property of \textbf{Rupam Ghosh} and \textbf{Studiva\textsuperscript{TM}}.\\[0.6em]
\textbf{Terms of Use \& Attribution Rules:}
\begin{itemize}[leftmargin=*, nosep]
    \item \textbf{Personal Educational Use:} Permitted freely for individual students and educators preparing for competitive examinations.
    \item \textbf{Mandatory Attribution:} Any reproduction, republication, hosting, or sharing (digital or physical) must prominently credit \textbf{Rupam Ghosh \& Studiva\textsuperscript{TM}} and include an active, direct link to \texttt{https://www.studiva.co.in}.
    \item \textbf{Commercial Restrictions:} Commercial resale, sublicensing, course bundling, or inclusion in paid test series without prior written permission is strictly prohibited.
    \item \textbf{Anti-Tampering Notice:} Removing, altering, or obscuring copyright notices, page watermarks, or website links from this document is strictly prohibited.
\end{itemize}
\vspace{0.6em}
\textbf{Fair Use Disclaimer:}\\
The underlying examination questions originate from the Graduate Aptitude Test in Engineering (GATE), conducted by the Indian Institutes of Technology (IITs) and Indian Institute of Science (IISc). Question texts are reproduced for non-commercial educational study under Fair Use doctrines. No copyright is claimed over the original GATE examination question texts.
\end{tcolorbox}

\vspace{1.2em}
\begin{center}
\small\itshape
Practice Recommendation: Use an opaque sheet of paper or ruler to cover the right-hand column during mock attempts. Tick your chosen answers directly in the check boxes on the left, then slide the cover to check the solution and hint.
\end{center}

\newpage
% ==================== PAGE 2 ONWARDS: QUESTION BANK ====================

\noindent
\begin{tcolorbox}[
    blanker,
    sidebyside,
    sidebyside align=top,
    lefthand width=0.55\textwidth,
    top=4pt, bottom=4pt, left=0pt, right=0pt,
    middle=16pt,
    underlay={
        \draw[dashed, black!80, line width=0.8pt, dash pattern=on 4pt off 3pt]
            ([xshift=0.55\textwidth+8pt]frame.north west) --
            ([xshift=0.55\textwidth+8pt]frame.south west);
    }
]
\textbf{Q.1} \quad {\footnotesize\textsc{[MCQ $\cdot$ 2 Mark(s) $\cdot$ Neg: 0.66]}} \hfill {\footnotesize\textbf{GATE 2018 CS} \quad \tiny\texttt{2018\_CS\_GA\_Q30}}
\vspace{2pt}\hrule height 0.3pt\vspace{4pt}

Let $G$ be a simple undirected graph. Let $T_D$ be a depth first search tree of $G$. Let $T_B$ be a breadth first search tree of $G$. Consider the following statements:
\begin{enumerate}[label=\Roman*.]
    \item For every edge $(u, v)$ of $G$, if $u$ is at depth $i$ and $v$ is at depth $j$ in $T_D$, then $|i - j| \ge 1$.
    \item For every edge $(u, v)$ of $G$, if $u$ is at depth $i$ and $v$ is at depth $j$ in $T_B$, then $|i - j| = 1$.
\end{enumerate}
Which of the statements above is/are necessarily TRUE?

\par\vspace{6pt}\noindent
\begin{tabular}{@{}p{0.48\linewidth}@{\hspace{0.04\linewidth}}p{0.48\linewidth}@{}}
$\square$\ \textbf{(A)} I only & $\square$\ \textbf{(B)} II only \\
\addlinespace[3pt]
$\square$\ \textbf{(C)} Both I and II & $\square$\ \textbf{(D)} Neither I nor II
\end{tabular}

\tcblower

\textbf{Official Key:} \textbf{\large A} \hfill \textit{\small [Medium]}
\vspace{2pt}\hrule height 0.3pt\vspace{3pt}
\textbf{Topic:} Algorithms \\
\textbf{Subtopic:} \textit{Graph traversals, minimum spanning trees, shortest paths}
\vspace{2pt}\hrule height 0.3pt\vspace{3pt}
\textbf{Guide / Hint:} \\
{\small In a DFS tree of an undirected graph, every non-tree edge is a back edge connecting an ancestor and a descendant, so forward and cross edges do not exist; thus statement I is necessarily true. In a BFS tree, cross edges can connect vertices at the same level ($|i - j| = 0$), so statement II is false.}

\end{tcolorbox}
\vspace{6pt}\hrule height 0.35pt\vspace{6pt}
\noindent
\begin{tcolorbox}[
    blanker,
    sidebyside,
    sidebyside align=top,
    lefthand width=0.55\textwidth,
    top=4pt, bottom=4pt, left=0pt, right=0pt,
    middle=16pt,
    underlay={
        \draw[dashed, black!80, line width=0.8pt, dash pattern=on 4pt off 3pt]
            ([xshift=0.55\textwidth+8pt]frame.north west) --
            ([xshift=0.55\textwidth+8pt]frame.south west);
    }
]
\textbf{Q.2} \quad {\footnotesize\textsc{[MCQ $\cdot$ 2 Mark(s) $\cdot$ Neg: 0.66]}} \hfill {\footnotesize\textbf{GATE 2018 CS} \quad \tiny\texttt{2018\_CS\_GA\_Q31}}
\vspace{2pt}\hrule height 0.3pt\vspace{4pt}

Assume that multiplying a matrix $G_1$ of dimension $p \times q$ with another matrix $G_2$ of dimension $q \times r$ requires $pqr$ scalar multiplications. Computing the product of five matrices $F_1, F_2, F_3, F_4$ and $F_5$ of dimensions $2 \times 25$, $25 \times 3$, $3 \times 16$, $16 \times 1$ and $1 \times 10$, respectively, using the minimum number of scalar multiplications requires first computing

\par\vspace{6pt}\noindent
\begin{tabular}{@{}p{0.48\linewidth}@{\hspace{0.04\linewidth}}p{0.48\linewidth}@{}}
$\square$\ \textbf{(A)} $F_1 F_2$ and $F_3 F_4$ & $\square$\ \textbf{(B)} $F_2 F_3$ and $F_3 F_4$ \\
\addlinespace[3pt]
$\square$\ \textbf{(C)} $F_2 F_3$ and $F_4 F_5$ & $\square$\ \textbf{(D)} $F_1 F_2$ and $F_4 F_5$
\end{tabular}

\tcblower

\textbf{Official Key:} \textbf{\large C} \hfill \textit{\small [Medium]}
\vspace{2pt}\hrule height 0.3pt\vspace{3pt}
\textbf{Topic:} Algorithms \\
\textbf{Subtopic:} \textit{Algorithm design techniques: greedy, dynamic programming and divide-and-conquer}
\vspace{2pt}\hrule height 0.3pt\vspace{3pt}
\textbf{Guide / Hint:} \\
{\small Using dynamic programming for matrix chain multiplication with dimensions $d = \langle 2, 25, 3, 16, 1, 10 \rangle$: calculating minimum multiplication costs reveals optimal parenthesization as $((F_1(F_2 F_3))(F_4 F_5))$, where the intermediate pairs computed first are $F_2 F_3$ and $F_4 F_5$.}

\end{tcolorbox}
\vspace{6pt}\hrule height 0.35pt\vspace{6pt}
\noindent
\begin{tcolorbox}[
    blanker,
    sidebyside,
    sidebyside align=top,
    lefthand width=0.55\textwidth,
    top=4pt, bottom=4pt, left=0pt, right=0pt,
    middle=16pt,
    underlay={
        \draw[dashed, black!80, line width=0.8pt, dash pattern=on 4pt off 3pt]
            ([xshift=0.55\textwidth+8pt]frame.north west) --
            ([xshift=0.55\textwidth+8pt]frame.south west);
    }
]
\textbf{Q.3} \quad {\footnotesize\textsc{[MCQ $\cdot$ 2 Mark(s) $\cdot$ Neg: 0.66]}} \hfill {\footnotesize\textbf{GATE 2018 CS} \quad \tiny\texttt{2018\_CS\_GA\_Q32}}
\vspace{2pt}\hrule height 0.3pt\vspace{4pt}

Consider the following C code. Assume that \texttt{unsigned long int} type length is 64 bits.
\begin{lstlisting}[language=C]
unsigned long int fun(unsigned long int n) {
    unsigned long int i, j = 0, sum = 0;
    while (n > 1) {
        j++;
        n /= 2;
    }
    for (i = 0; j > 1; i++) {
        sum += (j % 2);
        j /= 2;
    }
    return sum;
}
\end{lstlisting}
The value returned when we call \texttt{fun(pow(2, 40))} is

\par\vspace{6pt}\noindent
\begin{tabular}{@{}p{0.48\linewidth}@{\hspace{0.04\linewidth}}p{0.48\linewidth}@{}}
$\square$\ \textbf{(A)} 4 & $\square$\ \textbf{(B)} 5 \\
\addlinespace[3pt]
$\square$\ \textbf{(C)} 6 & $\square$\ \textbf{(D)} 40
\end{tabular}

\tcblower

\textbf{Official Key:} \textbf{\large B} \hfill \textit{\small [Medium]}
\vspace{2pt}\hrule height 0.3pt\vspace{3pt}
\textbf{Topic:} Algorithms \\
\textbf{Subtopic:} \textit{Asymptotic worst case time and space complexity}
\vspace{2pt}\hrule height 0.3pt\vspace{3pt}
\textbf{Guide / Hint:} \\
{\small For $n = 2^{40}$, the first loop `n /= 2` executes until $n = 1$, which takes $\log_2(2^{40}) = 40$ steps, giving $j = 40$. The second loop starts with $j = 40$ and halves $j$ until $j \le 1$: values taken by $j$ are $40 \to 20 \to 10 \to 5 \to 2 \to 1$, incrementing `sum` exactly 5 times.}

\end{tcolorbox}
\vspace{6pt}\hrule height 0.35pt\vspace{6pt}
\noindent
\begin{tcolorbox}[
    blanker,
    sidebyside,
    sidebyside align=top,
    lefthand width=0.55\textwidth,
    top=4pt, bottom=4pt, left=0pt, right=0pt,
    middle=16pt,
    underlay={
        \draw[dashed, black!80, line width=0.8pt, dash pattern=on 4pt off 3pt]
            ([xshift=0.55\textwidth+8pt]frame.north west) --
            ([xshift=0.55\textwidth+8pt]frame.south west);
    }
]
\textbf{Q.4} \quad {\footnotesize\textsc{[NAT $\cdot$ 2 Mark(s) $\cdot$ No Neg]}} \hfill {\footnotesize\textbf{GATE 2018 CS} \quad \tiny\texttt{2018\_CS\_GA\_Q48}}
\vspace{2pt}\hrule height 0.3pt\vspace{4pt}

Consider the following items with weights and values:
\begin{center}
\begin{tabular}{ccc}
\textbf{Item} & \textbf{Weight} & \textbf{Value} \\
1 & 10 & 60 \\
2 & 7 & 28 \\
3 & 4 & 20 \\
4 & 2 & 24 \\
\end{tabular}
\end{center}
A knapsack has capacity 11. A greedy heuristic sorts items in decreasing order of value-to-weight ratio and picks each item that fits. Let $V_{\text{opt}}$ and $V_{\text{heuristic}}$ denote the optimal knapsack value and heuristic value, respectively. The value of $V_{\text{opt}} - V_{\text{heuristic}}$ is \_\_\_\_\_\_\_\_.

\par\vspace{6pt}\noindent\textbf{Your Answer:} \fbox{\rule{0pt}{1.3ex}\hspace{3cm}}

\tcblower

\textbf{Official Key:} \textbf{\large 16 to 16} \hfill \textit{\small [Medium]}
\vspace{2pt}\hrule height 0.3pt\vspace{3pt}
\textbf{Topic:} Algorithms \\
\textbf{Subtopic:} \textit{Algorithm design techniques: greedy, dynamic programming and divide-and-conquer}
\vspace{2pt}\hrule height 0.3pt\vspace{3pt}
\textbf{Guide / Hint:} \\
{\small Value-to-weight ratios: Item 4 ($24/2 = 12$), Item 1 ($60/10 = 6$), Item 3 ($20/4 = 5$), Item 2 ($28/7 = 4$). The heuristic sorts by ratio and greedily takes Item 4 (wt 2, remaining 9); Item 1 does not fit ($10 > 9$); Item 3 fits (wt 4, remaining 5); Item 2 does not fit ($7 > 5$), yielding $V_{\text{heuristic}} = 24 + 20 = 44$. The optimal 0/1 knapsack selection picks Item 1 (wt 10, val 60). Thus $V_{\text{opt}} - V_{\text{heuristic}} = 60 - 44 = 16$.}

\end{tcolorbox}
\vspace{6pt}\hrule height 0.35pt\vspace{6pt}
\noindent
\begin{tcolorbox}[
    blanker,
    sidebyside,
    sidebyside align=top,
    lefthand width=0.55\textwidth,
    top=4pt, bottom=4pt, left=0pt, right=0pt,
    middle=16pt,
    underlay={
        \draw[dashed, black!80, line width=0.8pt, dash pattern=on 4pt off 3pt]
            ([xshift=0.55\textwidth+8pt]frame.north west) --
            ([xshift=0.55\textwidth+8pt]frame.south west);
    }
]
\textbf{Q.5} \quad {\footnotesize\textsc{[MCQ $\cdot$ 1 Mark(s) $\cdot$ Neg: 0.33]}} \hfill {\footnotesize\textbf{GATE 2022 CS} \quad \tiny\texttt{2022\_CS\_CS\_Q11}}
\vspace{2pt}\hrule height 0.3pt\vspace{4pt}

Which one of the following statements is TRUE for all positive functions $f(n)$?

\par\vspace{6pt}\noindent
\begin{tabular}{@{}p{0.98\linewidth}@{}}
$\square$\ \textbf{(A)} $f(n^2) = \Theta((f(n))^2)$, when $f(n)$ is a polynomial \\
\addlinespace[3pt]
$\square$\ \textbf{(B)} $f(n^2) = o((f(n))^2)$ \\
\addlinespace[3pt]
$\square$\ \textbf{(C)} $f(n^2) = \mathcal{O}((f(n))^2)$, when $f(n)$ is an exponential function \\
\addlinespace[3pt]
$\square$\ \textbf{(D)} $f(n^2) = \Omega((f(n))^2)$
\end{tabular}

\tcblower

\textbf{Official Key:} \textbf{\large A} \hfill \textit{\small [Medium]}
\vspace{2pt}\hrule height 0.3pt\vspace{3pt}
\textbf{Topic:} Algorithms \\
\textbf{Subtopic:} \textit{Asymptotic worst case time and space complexity}
\vspace{2pt}\hrule height 0.3pt\vspace{3pt}
\textbf{Guide / Hint:} \\
{\small If $f(n)$ is a polynomial of degree $k$ ($f(n) = \Theta(n^k)$), then $f(n^2) = \Theta((n^2)^k) = \Theta(n^{2k})$, and $(f(n))^2 = (\Theta(n^k))^2 = \Theta(n^{2k})$. Since both functions grow at identical polynomial rates, $f(n^2) = \Theta((f(n))^2)$ holds.}

\end{tcolorbox}
\vspace{6pt}\hrule height 0.35pt\vspace{6pt}
\noindent
\begin{tcolorbox}[
    blanker,
    sidebyside,
    sidebyside align=top,
    lefthand width=0.55\textwidth,
    top=4pt, bottom=4pt, left=0pt, right=0pt,
    middle=16pt,
    underlay={
        \draw[dashed, black!80, line width=0.8pt, dash pattern=on 4pt off 3pt]
            ([xshift=0.55\textwidth+8pt]frame.north west) --
            ([xshift=0.55\textwidth+8pt]frame.south west);
    }
]
\textbf{Q.6} \quad {\footnotesize\textsc{[MCQ $\cdot$ 1 Mark(s) $\cdot$ Neg: 0.33]}} \hfill {\footnotesize\textbf{GATE 2022 CS} \quad \tiny\texttt{2022\_CS\_CS\_Q16}}
\vspace{2pt}\hrule height 0.3pt\vspace{4pt}

Suppose we are given $n$ keys, $m$ hash table slots, and two simple uniform hash functions $h_1$ and $h_2$. Further suppose our hashing scheme uses $h_1$ for the odd keys and $h_2$ for the even keys. What is the expected number of keys in a slot?

\par\vspace{6pt}\noindent
\begin{tabular}{@{}p{0.48\linewidth}@{\hspace{0.04\linewidth}}p{0.48\linewidth}@{}}
$\square$\ \textbf{(A)} $m/n$ & $\square$\ \textbf{(B)} $n/m$ \\
\addlinespace[3pt]
$\square$\ \textbf{(C)} $2n/m$ & $\square$\ \textbf{(D)} $n / (2m)$
\end{tabular}

\tcblower

\textbf{Official Key:} \textbf{\large B} \hfill \textit{\small [Easy]}
\vspace{2pt}\hrule height 0.3pt\vspace{3pt}
\textbf{Topic:} Algorithms \\
\textbf{Subtopic:} \textit{Searching, sorting, hashing}
\vspace{2pt}\hrule height 0.3pt\vspace{3pt}
\textbf{Guide / Hint:} \\
{\small Let $n_1$ be the number of odd keys and $n_2$ be the number of even keys ($n_1 + n_2 = n$). Under simple uniform hashing, each key independently hashes to any slot with probability $1/m$. By linearity of expectation, the expected number of keys in any slot is $\frac{n_1}{m} + \frac{n_2}{m} = \frac{n}{m}$.}

\end{tcolorbox}
\vspace{6pt}\hrule height 0.35pt\vspace{6pt}
\noindent
\begin{tcolorbox}[
    blanker,
    sidebyside,
    sidebyside align=top,
    lefthand width=0.55\textwidth,
    top=4pt, bottom=4pt, left=0pt, right=0pt,
    middle=16pt,
    underlay={
        \draw[dashed, black!80, line width=0.8pt, dash pattern=on 4pt off 3pt]
            ([xshift=0.55\textwidth+8pt]frame.north west) --
            ([xshift=0.55\textwidth+8pt]frame.south west);
    }
]
\textbf{Q.7} \quad {\footnotesize\textsc{[MCQ $\cdot$ 2 Mark(s) $\cdot$ Neg: 0.66]}} \hfill {\footnotesize\textbf{GATE 2022 CS} \quad \tiny\texttt{2022\_CS\_CS\_Q37}}
\vspace{2pt}\hrule height 0.3pt\vspace{4pt}

Consider a simple undirected unweighted graph with at least three vertices. If $A$ is the adjacency matrix of the graph, then the number of 3-cycles in the graph is given by the trace of

\par\vspace{6pt}\noindent
\begin{tabular}{@{}p{0.48\linewidth}@{\hspace{0.04\linewidth}}p{0.48\linewidth}@{}}
$\square$\ \textbf{(A)} $A^3$ & $\square$\ \textbf{(B)} $A^3$ divided by 2 \\
\addlinespace[3pt]
$\square$\ \textbf{(C)} $A^3$ divided by 3 & $\square$\ \textbf{(D)} $A^3$ divided by 6
\end{tabular}

\tcblower

\textbf{Official Key:} \textbf{\large D} \hfill \textit{\small [Medium]}
\vspace{2pt}\hrule height 0.3pt\vspace{3pt}
\textbf{Topic:} Algorithms \\
\textbf{Subtopic:} \textit{Graph traversals, minimum spanning trees, shortest paths}
\vspace{2pt}\hrule height 0.3pt\vspace{3pt}
\textbf{Guide / Hint:} \\
{\small The diagonal entry $(A^3)_{ii}$ counts the number of closed walks of length 3 starting and ending at vertex $i$, so $\text{tr}(A^3) = \sum (A^3)_{ii}$ counts all closed walks of length 3 in $G$. In a simple undirected graph, each triangle contains 3 vertices and can be traversed in 2 directions, counting each 3-cycle $3 \times 2 = 6$ times. Thus, the number of 3-cycles is $\frac{\text{tr}(A^3)}{6}$.}

\end{tcolorbox}
\vspace{6pt}\hrule height 0.35pt\vspace{6pt}
\noindent
\begin{tcolorbox}[
    blanker,
    sidebyside,
    sidebyside align=top,
    lefthand width=0.55\textwidth,
    top=4pt, bottom=4pt, left=0pt, right=0pt,
    middle=16pt,
    underlay={
        \draw[dashed, black!80, line width=0.8pt, dash pattern=on 4pt off 3pt]
            ([xshift=0.55\textwidth+8pt]frame.north west) --
            ([xshift=0.55\textwidth+8pt]frame.south west);
    }
]
\textbf{Q.8} \quad {\footnotesize\textsc{[MSQ $\cdot$ 2 Mark(s) $\cdot$ No Neg]}} \hfill {\footnotesize\textbf{GATE 2022 CS} \quad \tiny\texttt{2022\_CS\_CS\_Q49}}
\vspace{2pt}\hrule height 0.3pt\vspace{4pt}

Consider a simple undirected weighted graph $G$, all of whose edge weights are distinct. Which of the following statements about the minimum spanning trees of $G$ is/are TRUE?

\par\vspace{6pt}\noindent
\begin{tabular}{@{}p{0.98\linewidth}@{}}
$\square$\ \textbf{(A)} The edge with the second smallest weight is always part of any minimum spanning tree of $G$. \\
\addlinespace[3pt]
$\square$\ \textbf{(B)} One or both of the edges with the third smallest and the fourth smallest weights are part of any minimum spanning tree of $G$. \\
\addlinespace[3pt]
$\square$\ \textbf{(C)} Suppose $S \subset V$ be such that $S \neq \emptyset$ and $S \neq V$. Consider the edge with the minimum weight such that one of its vertices is in $S$ and the other in $V \setminus S$. Such an edge will always be part of any minimum spanning tree of $G$. \\
\addlinespace[3pt]
$\square$\ \textbf{(D)} $G$ can have multiple minimum spanning trees.
\end{tabular}

\tcblower

\textbf{Official Key:} \textbf{\large A,B,C} \hfill \textit{\small [Medium]}
\vspace{2pt}\hrule height 0.3pt\vspace{3pt}
\textbf{Topic:} Algorithms \\
\textbf{Subtopic:} \textit{Graph traversals, minimum spanning trees, shortest paths}
\vspace{2pt}\hrule height 0.3pt\vspace{3pt}
\textbf{Guide / Hint:} \\
{\small The two edges with the smallest weights can never form a cycle, so both belong to the MST (A is TRUE). The three smallest edges can at most form one cycle (triangle), so at least one of the 3rd and 4th smallest edges must be part of the MST (B is TRUE). By the Cut Property, the minimum weight edge across any non-trivial cut $(S, V \setminus S)$ belongs to the MST (C is TRUE). Distinct weights guarantee a unique MST (D is FALSE).}

\end{tcolorbox}
\vspace{6pt}\hrule height 0.35pt\vspace{6pt}
\noindent
\begin{tcolorbox}[
    blanker,
    sidebyside,
    sidebyside align=top,
    lefthand width=0.55\textwidth,
    top=4pt, bottom=4pt, left=0pt, right=0pt,
    middle=16pt,
    underlay={
        \draw[dashed, black!80, line width=0.8pt, dash pattern=on 4pt off 3pt]
            ([xshift=0.55\textwidth+8pt]frame.north west) --
            ([xshift=0.55\textwidth+8pt]frame.south west);
    }
]
\textbf{Q.9} \quad {\footnotesize\textsc{[MSQ $\cdot$ 2 Mark(s) $\cdot$ No Neg]}} \hfill {\footnotesize\textbf{GATE 2022 CS} \quad \tiny\texttt{2022\_CS\_CS\_Q51}}
\vspace{2pt}\hrule height 0.3pt\vspace{4pt}

Consider the following recurrence:
\[ f(1) = 1; \]
\[ f(2n) = 2f(n) - 1, \quad \text{for } n \ge 1; \]
\[ f(2n+1) = 2f(n) + 1, \quad \text{for } n \ge 1. \]
Then, which of the following statements is/are TRUE?

\par\vspace{6pt}\noindent
\begin{tabular}{@{}p{0.48\linewidth}@{\hspace{0.04\linewidth}}p{0.48\linewidth}@{}}
$\square$\ \textbf{(A)} $f(2^n - 1) = 2^n - 1$ & $\square$\ \textbf{(B)} $f(2^n) = 1$ \\
\addlinespace[3pt]
$\square$\ \textbf{(C)} $f(5 \cdot 2^n) = 2^{n+1} + 1$ & $\square$\ \textbf{(D)} $f(2^n + 1) = 2^n + 1$
\end{tabular}

\tcblower

\textbf{Official Key:} \textbf{\large A,B,C} \hfill \textit{\small [Medium]}
\vspace{2pt}\hrule height 0.3pt\vspace{3pt}
\textbf{Topic:} Algorithms \\
\textbf{Subtopic:} \textit{Asymptotic worst case time and space complexity}
\vspace{2pt}\hrule height 0.3pt\vspace{3pt}
\textbf{Guide / Hint:} \\
{\small Computing small values: $f(1)=1, f(2)=1, f(3)=3, f(4)=1, f(5)=5, f(6)=5, f(7)=7, f(8)=1$. In binary, $f(n)$ corresponds to cyclic left shift of $n$. For $2^n - 1$ (all 1s in binary), $f(2^n - 1) = 2^n - 1$ (A is TRUE). For powers of 2, $f(2^n) = 1$ (B is TRUE). For $5 \cdot 2^n$, binary representation is $10100\dots$, shifting gives $100\dots01_2 = 2^{n+1} + 1$ (C is TRUE).}

\end{tcolorbox}
\vspace{6pt}\hrule height 0.35pt\vspace{6pt}
\noindent
\begin{tcolorbox}[
    blanker,
    sidebyside,
    sidebyside align=top,
    lefthand width=0.55\textwidth,
    top=4pt, bottom=4pt, left=0pt, right=0pt,
    middle=16pt,
    underlay={
        \draw[dashed, black!80, line width=0.8pt, dash pattern=on 4pt off 3pt]
            ([xshift=0.55\textwidth+8pt]frame.north west) --
            ([xshift=0.55\textwidth+8pt]frame.south west);
    }
]
\textbf{Q.10} \quad {\footnotesize\textsc{[NAT $\cdot$ 2 Mark(s) $\cdot$ No Neg]}} \hfill {\footnotesize\textbf{GATE 2022 CS} \quad \tiny\texttt{2022\_CS\_CS\_Q58}}
\vspace{2pt}\hrule height 0.3pt\vspace{4pt}

Let $G = (V, E)$ be a directed graph, where $V = \{1, 2, 3, 4, 5\}$ is the set of vertices and $E$ is the set of directed edges, as defined by the adjacency matrix $A[i][j] = 1$ if $1 \le j \le i \le 5$, and 0 otherwise. A directed spanning tree of $G$ rooted at $r \in V$ is a subgraph $T$ whose undirected version is a tree and $T$ contains a directed path from $r$ to every other vertex. The number of such directed spanning trees rooted at vertex 5 is \_\_\_\_\_\_\_\_\_\_\_\_\. (Answer in integer)

\par\vspace{6pt}\noindent\textbf{Your Answer:} \fbox{\rule{0pt}{1.3ex}\hspace{3cm}}

\tcblower

\textbf{Official Key:} \textbf{\large 24 to 24} \hfill \textit{\small [Hard]}
\vspace{2pt}\hrule height 0.3pt\vspace{3pt}
\textbf{Topic:} Algorithms \\
\textbf{Subtopic:} \textit{Graph traversals, minimum spanning trees, shortest paths}
\vspace{2pt}\hrule height 0.3pt\vspace{3pt}
\textbf{Guide / Hint:} \\
{\small In a directed spanning tree rooted at 5 with edges $A[i][j]=1$ for $j \le i$, every node $v \in \{1, 2, 3, 4\}$ must choose exactly one parent $u$ with $u > v$. Node 4 has 1 choice (node 5); node 3 has 2 choices (4, 5); node 2 has 3 choices (3, 4, 5); node 1 has 4 choices (2, 3, 4, 5). Total spanning trees $= 1 \times 2 \times 3 \times 4 = 24$.}

\end{tcolorbox}
\vspace{6pt}\hrule height 0.35pt\vspace{6pt}
\noindent
\begin{tcolorbox}[
    blanker,
    sidebyside,
    sidebyside align=top,
    lefthand width=0.55\textwidth,
    top=4pt, bottom=4pt, left=0pt, right=0pt,
    middle=16pt,
    underlay={
        \draw[dashed, black!80, line width=0.8pt, dash pattern=on 4pt off 3pt]
            ([xshift=0.55\textwidth+8pt]frame.north west) --
            ([xshift=0.55\textwidth+8pt]frame.south west);
    }
]
\textbf{Q.11} \quad {\footnotesize\textsc{[MCQ $\cdot$ 1 Mark(s) $\cdot$ Neg: 0.33]}} \hfill {\footnotesize\textbf{GATE 2023 CS} \quad \tiny\texttt{2023\_CS\_CS\_Q20}}
\vspace{2pt}\hrule height 0.3pt\vspace{4pt}

An algorithm has to store several keys generated by an adversary in a hash table. The adversary is malicious who tries to maximize the number of collisions. Let $k$ be the number of keys, $m$ be the number of slots in the hash table, and $k > m$.
Which one of the following is the best hashing strategy to counteract the adversary?

\par\vspace{6pt}\noindent
\begin{tabular}{@{}p{0.98\linewidth}@{}}
$\square$\ \textbf{(A)} Division method, i.e., use the hash function $h(k) = k \bmod m$. \\
\addlinespace[3pt]
$\square$\ \textbf{(B)} Multiplication method, i.e., use the hash function $h(k) = \lfloor m(kA - \lfloor kA \rfloor)\rfloor$, where $A$ is a carefully chosen constant. \\
\addlinespace[3pt]
$\square$\ \textbf{(C)} Universal hashing method. \\
\addlinespace[3pt]
$\square$\ \textbf{(D)} If $k$ is a prime number, use Division method. Otherwise, use Multiplication method.
\end{tabular}

\tcblower

\textbf{Official Key:} \textbf{\large C} \hfill \textit{\small [Easy]}
\vspace{2pt}\hrule height 0.3pt\vspace{3pt}
\textbf{Topic:} Algorithms \\
\textbf{Subtopic:} \textit{Searching, sorting, hashing}
\vspace{2pt}\hrule height 0.3pt\vspace{3pt}
\textbf{Guide / Hint:} \\
{\small In universal hashing, the hash function is chosen uniformly at random from a carefully designed universal collection of hash functions independently of the keys. Because the adversary cannot predict which function will be selected at runtime, any fixed choice of malicious keys cannot guarantee worst-case collision behavior.}

\end{tcolorbox}
\vspace{6pt}\hrule height 0.35pt\vspace{6pt}
\noindent
\begin{tcolorbox}[
    blanker,
    sidebyside,
    sidebyside align=top,
    lefthand width=0.55\textwidth,
    top=4pt, bottom=4pt, left=0pt, right=0pt,
    middle=16pt,
    underlay={
        \draw[dashed, black!80, line width=0.8pt, dash pattern=on 4pt off 3pt]
            ([xshift=0.55\textwidth+8pt]frame.north west) --
            ([xshift=0.55\textwidth+8pt]frame.south west);
    }
]
\textbf{Q.12} \quad {\footnotesize\textsc{[MSQ $\cdot$ 1 Mark(s) $\cdot$ No Neg]}} \hfill {\footnotesize\textbf{GATE 2023 CS} \quad \tiny\texttt{2023\_CS\_CS\_Q29}}
\vspace{2pt}\hrule height 0.3pt\vspace{4pt}

Let $f$ and $g$ be functions of natural numbers given by $f(n) = n$ and $g(n) = n^2$.
Which of the following statements is/are TRUE?

\par\vspace{6pt}\noindent
\begin{tabular}{@{}p{0.48\linewidth}@{\hspace{0.04\linewidth}}p{0.48\linewidth}@{}}
$\square$\ \textbf{(A)} $f \in \mathcal{O}(g)$ & $\square$\ \textbf{(B)} $f \in \Omega(g)$ \\
\addlinespace[3pt]
$\square$\ \textbf{(C)} $f \in o(g)$ & $\square$\ \textbf{(D)} $f \in \Theta(g)$
\end{tabular}

\tcblower

\textbf{Official Key:} \textbf{\large A, C} \hfill \textit{\small [Easy]}
\vspace{2pt}\hrule height 0.3pt\vspace{3pt}
\textbf{Topic:} Algorithms \\
\textbf{Subtopic:} \textit{Asymptotic worst case time and space complexity}
\vspace{2pt}\hrule height 0.3pt\vspace{3pt}
\textbf{Guide / Hint:} \\
{\small For $f(n) = n$ and $g(n) = n^2$, we evaluate $\lim_{n \to \infty} \frac{f(n)}{g(n)} = \lim_{n \to \infty} \frac{1}{n} = 0$. By definition, if the limit is 0, $f(n) \in o(g(n))$, which strictly implies $f(n) \in \mathcal{O}(g(n))$. Hence, both A and C are TRUE.}

\end{tcolorbox}
\vspace{6pt}\hrule height 0.35pt\vspace{6pt}
\noindent
\begin{tcolorbox}[
    blanker,
    sidebyside,
    sidebyside align=top,
    lefthand width=0.55\textwidth,
    top=4pt, bottom=4pt, left=0pt, right=0pt,
    middle=16pt,
    underlay={
        \draw[dashed, black!80, line width=0.8pt, dash pattern=on 4pt off 3pt]
            ([xshift=0.55\textwidth+8pt]frame.north west) --
            ([xshift=0.55\textwidth+8pt]frame.south west);
    }
]
\textbf{Q.13} \quad {\footnotesize\textsc{[MSQ $\cdot$ 2 Mark(s) $\cdot$ No Neg]}} \hfill {\footnotesize\textbf{GATE 2023 CS} \quad \tiny\texttt{2023\_CS\_CS\_Q54}}
\vspace{2pt}\hrule height 0.3pt\vspace{4pt}

Consider functions Function 1 and Function 2 expressed in pseudocode as follows:
\begin{lstlisting}
// Function 1
while n > 1 do
    for i = 1 to n do
        x = x + 1;
    end for
    n = floor(n / 2);
end while
\par\medskip
// Function 2
for i = 1 to 100 * n do
    x = x + 1;
end for
\end{lstlisting}
Let $f_1(n)$ and $f_2(n)$ denote the number of times the statement ``$x = x + 1$'' is executed in Function 1 and Function 2, respectively.
Which of the following statements is/are TRUE?

\par\vspace{6pt}\noindent
\begin{tabular}{@{}p{0.48\linewidth}@{\hspace{0.04\linewidth}}p{0.48\linewidth}@{}}
$\square$\ \textbf{(A)} $f_1(n) \in \Theta(f_2(n))$ & $\square$\ \textbf{(B)} $f_1(n) \in o(f_2(n))$ \\
\addlinespace[3pt]
$\square$\ \textbf{(C)} $f_1(n) \in \omega(f_2(n))$ & $\square$\ \textbf{(D)} $f_1(n) \in \mathcal{O}(n)$
\end{tabular}

\tcblower

\textbf{Official Key:} \textbf{\large A, D} \hfill \textit{\small [Easy]}
\vspace{2pt}\hrule height 0.3pt\vspace{3pt}
\textbf{Topic:} Algorithms \\
\textbf{Subtopic:} \textit{Asymptotic worst case time and space complexity}
\vspace{2pt}\hrule height 0.3pt\vspace{3pt}
\textbf{Guide / Hint:} \\
{\small In Function 1, the inner loop executes $n + \lfloor n/2 \rfloor + \lfloor n/4 \rfloor + \dots \approx 2n$ times, which gives $f_1(n) = \Theta(n)$. In Function 2, the loop runs $100n$ times, which gives $f_2(n) = \Theta(n)$. Because both functions have exact linear asymptotic growth, $f_1(n) \in \Theta(f_2(n))$ (A is TRUE) and $f_1(n) \in \mathcal{O}(n)$ (D is TRUE).}

\end{tcolorbox}
\vspace{6pt}\hrule height 0.35pt\vspace{6pt}
\noindent
\begin{tcolorbox}[
    blanker,
    sidebyside,
    sidebyside align=top,
    lefthand width=0.55\textwidth,
    top=4pt, bottom=4pt, left=0pt, right=0pt,
    middle=16pt,
    underlay={
        \draw[dashed, black!80, line width=0.8pt, dash pattern=on 4pt off 3pt]
            ([xshift=0.55\textwidth+8pt]frame.north west) --
            ([xshift=0.55\textwidth+8pt]frame.south west);
    }
]
\textbf{Q.14} \quad {\footnotesize\textsc{[NAT $\cdot$ 2 Mark(s) $\cdot$ No Neg]}} \hfill {\footnotesize\textbf{GATE 2023 CS} \quad \tiny\texttt{2023\_CS\_CS\_Q56}}
\vspace{2pt}\hrule height 0.3pt\vspace{4pt}

Let $U = \{1, 2, 3\}$. Let $2^U$ denote the powerset of $U$. Consider an undirected graph $G$ whose vertex set is $2^U$. For any $A, B \in 2^U$, $(A, B)$ is an edge in $G$ if and only if (i) $A \neq B$, and (ii) either $A \subsetneq B$ or $B \subsetneq A$. For any vertex $A$ in $G$, the set of all possible orderings in which the vertices of $G$ can be visited in a Breadth First Search (BFS) starting from $A$ is denoted by $B(A)$.
If $\emptyset$ denotes the empty set, then the cardinality of $B(\emptyset)$ is \_\_\_\_\_\_\_\_\. (Answer in integer)

\par\vspace{6pt}\noindent\textbf{Your Answer:} \fbox{\rule{0pt}{1.3ex}\hspace{3cm}}

\tcblower

\textbf{Official Key:} \textbf{\large 5040 to 5040} \hfill \textit{\small [Hard]}
\vspace{2pt}\hrule height 0.3pt\vspace{3pt}
\textbf{Topic:} Algorithms \\
\textbf{Subtopic:} \textit{Graph traversals, minimum spanning trees, shortest paths}
\vspace{2pt}\hrule height 0.3pt\vspace{3pt}
\textbf{Guide / Hint:} \\
{\small The vertex set has $2^3 = 8$ vertices. An edge connects $A$ and $B$ iff $A \subsetneq B$ or $B \subsetneq A$. For the empty set $\emptyset$, every non-empty subset $S \subseteq U$ satisfies $\emptyset \subsetneq S$, meaning all 7 non-empty subsets are direct neighbors of $\emptyset$ at BFS depth 1. A BFS starting at $\emptyset$ must visit $\emptyset$ first, and can then visit the 7 distance-1 neighbors in any of the $7! = 5040$ possible orderings.}

\end{tcolorbox}
\vspace{6pt}\hrule height 0.35pt\vspace{6pt}
\noindent
\begin{tcolorbox}[
    blanker,
    sidebyside,
    sidebyside align=top,
    lefthand width=0.55\textwidth,
    top=4pt, bottom=4pt, left=0pt, right=0pt,
    middle=16pt,
    underlay={
        \draw[dashed, black!80, line width=0.8pt, dash pattern=on 4pt off 3pt]
            ([xshift=0.55\textwidth+8pt]frame.north west) --
            ([xshift=0.55\textwidth+8pt]frame.south west);
    }
]
\textbf{Q.15} \quad {\footnotesize\textsc{[MCQ $\cdot$ 1 Mark(s) $\cdot$ Neg: 0.33]}} \hfill {\footnotesize\textbf{GATE 2024 CS1} \quad \tiny\texttt{2024\_CS1\_CS\_Q17}}
\vspace{2pt}\hrule height 0.3pt\vspace{4pt}

Given an integer array of size $N$, we want to check if the array is sorted (in either ascending or descending order). An algorithm solves this problem by making a single pass through the array and comparing each element of the array only with its adjacent elements. The worst-case time complexity of this algorithm is

\par\vspace{6pt}\noindent
\begin{tabular}{@{}p{0.98\linewidth}@{}}
$\square$\ \textbf{(A)} both $\mathcal{O}(N)$ and $\Omega(N)$ \\
\addlinespace[3pt]
$\square$\ \textbf{(B)} $\mathcal{O}(N)$ but not $\Omega(N)$ \\
\addlinespace[3pt]
$\square$\ \textbf{(C)} $\Omega(N)$ but not $\mathcal{O}(N)$ \\
\addlinespace[3pt]
$\square$\ \textbf{(D)} neither $\mathcal{O}(N)$ nor $\Omega(N)$
\end{tabular}

\tcblower

\textbf{Official Key:} \textbf{\large A} \hfill \textit{\small [Easy]}
\vspace{2pt}\hrule height 0.3pt\vspace{3pt}
\textbf{Topic:} Algorithms \\
\textbf{Subtopic:} \textit{Asymptotic worst case time and space complexity}
\vspace{2pt}\hrule height 0.3pt\vspace{3pt}
\textbf{Guide / Hint:} \\
{\small In the worst case (such as when the array is sorted or the disorder occurs at the very end), a single-pass algorithm that compares adjacent elements must traverse the entire array, performing $N-1$ comparisons. The worst-case runtime is asymptotically tight at $\Theta(N)$, meaning it is both $\mathcal{O}(N)$ and $\Omega(N)$.}

\end{tcolorbox}
\vspace{6pt}\hrule height 0.35pt\vspace{6pt}
\noindent
\begin{tcolorbox}[
    blanker,
    sidebyside,
    sidebyside align=top,
    lefthand width=0.55\textwidth,
    top=4pt, bottom=4pt, left=0pt, right=0pt,
    middle=16pt,
    underlay={
        \draw[dashed, black!80, line width=0.8pt, dash pattern=on 4pt off 3pt]
            ([xshift=0.55\textwidth+8pt]frame.north west) --
            ([xshift=0.55\textwidth+8pt]frame.south west);
    }
]
\textbf{Q.16} \quad {\footnotesize\textsc{[MCQ $\cdot$ 2 Mark(s) $\cdot$ Neg: 0.66]}} \hfill {\footnotesize\textbf{GATE 2024 CS1} \quad \tiny\texttt{2024\_CS1\_CS\_Q42}}
\vspace{2pt}\hrule height 0.3pt\vspace{4pt}

Consider the following recurrence relation:
\[ T(n) = \begin{cases} \sqrt{n} T(\sqrt{n}) + n & \text{for } n > 1, \\ 1 & \text{for } n = 1. \end{cases} \]
Which one of the following options is CORRECT?

\par\vspace{6pt}\noindent
\begin{tabular}{@{}p{0.98\linewidth}@{}}
$\square$\ \textbf{(A)} $T(n) = \Theta(n \log \log n)$ \\
\addlinespace[3pt]
$\square$\ \textbf{(B)} $T(n) = \Theta(n \log n)$ \\
\addlinespace[3pt]
$\square$\ \textbf{(C)} $T(n) = \Theta(n^2 \log n)$ \\
\addlinespace[3pt]
$\square$\ \textbf{(D)} $T(n) = \Theta(n^2 \log \log n)$
\end{tabular}

\tcblower

\textbf{Official Key:} \textbf{\large A} \hfill \textit{\small [Medium]}
\vspace{2pt}\hrule height 0.3pt\vspace{3pt}
\textbf{Topic:} Algorithms \\
\textbf{Subtopic:} \textit{Asymptotic worst case time and space complexity}
\vspace{2pt}\hrule height 0.3pt\vspace{3pt}
\textbf{Guide / Hint:} \\
{\small Dividing $T(n) = \sqrt{n} T(\sqrt{n}) + n$ by $n$ gives $\frac{T(n)}{n} = \frac{T(\sqrt{n})}{\sqrt{n}} + 1$. Defining $S(n) = \frac{T(n)}{n}$ yields $S(n) = S(\sqrt{n}) + 1$. Substituting $n = 2^{2^k}$ gives $S(2^{2^k}) = S(2^{2^{k-1}}) + 1 = \Theta(k) = \Theta(\log \log n)$. Therefore, $T(n) = n S(n) = \Theta(n \log \log n)$.}

\end{tcolorbox}
\vspace{6pt}\hrule height 0.35pt\vspace{6pt}
\noindent
\begin{tcolorbox}[
    blanker,
    sidebyside,
    sidebyside align=top,
    lefthand width=0.55\textwidth,
    top=4pt, bottom=4pt, left=0pt, right=0pt,
    middle=16pt,
    underlay={
        \draw[dashed, black!80, line width=0.8pt, dash pattern=on 4pt off 3pt]
            ([xshift=0.55\textwidth+8pt]frame.north west) --
            ([xshift=0.55\textwidth+8pt]frame.south west);
    }
]
\textbf{Q.17} \quad {\footnotesize\textsc{[MSQ $\cdot$ 2 Mark(s) $\cdot$ No Neg]}} \hfill {\footnotesize\textbf{GATE 2024 CS1} \quad \tiny\texttt{2024\_CS1\_CS\_Q45}}
\vspace{2pt}\hrule height 0.3pt\vspace{4pt}

Let $G$ be a directed graph and $T$ a depth first search (DFS) spanning tree in $G$ that is rooted at a vertex $v$. Suppose $T$ is also a breadth first search (BFS) tree in $G$, rooted at $v$. Which of the following statements is/are TRUE for every such graph $G$ and tree $T$?

\par\vspace{6pt}\noindent
\begin{tabular}{@{}p{0.98\linewidth}@{}}
$\square$\ \textbf{(A)} There are no back-edges in $G$ with respect to the tree $T$ \\
\addlinespace[3pt]
$\square$\ \textbf{(B)} There are no cross-edges in $G$ with respect to the tree $T$ \\
\addlinespace[3pt]
$\square$\ \textbf{(C)} There are no forward-edges in $G$ with respect to the tree $T$ \\
\addlinespace[3pt]
$\square$\ \textbf{(D)} The only edges in $G$ are the edges in $T$
\end{tabular}

\tcblower

\textbf{Official Key:} \textbf{\large C} \hfill \textit{\small [Hard]}
\vspace{2pt}\hrule height 0.3pt\vspace{3pt}
\textbf{Topic:} Algorithms \\
\textbf{Subtopic:} \textit{Graph traversals, minimum spanning trees, shortest paths}
\vspace{2pt}\hrule height 0.3pt\vspace{3pt}
\textbf{Guide / Hint:} \\
{\small In a BFS tree $T$, any edge $(u, w)$ in $G$ must connect vertices whose BFS levels differ by at most 1 (i.e., $\text{level}(w) \le \text{level}(u) + 1$). A forward-edge in a DFS tree goes from ancestor $u$ to descendant $w$ with $\text{level}(w) \ge \text{level}(u) + 2$. If such an edge existed in $G$, BFS would explore $w$ directly from $u$ at depth $\text{level}(u) + 1$, contradicting the structure of $T$. Hence, no forward-edges can exist in $G$ with respect to $T$ (C is TRUE).}

\end{tcolorbox}
\vspace{6pt}\hrule height 0.35pt\vspace{6pt}
\noindent
\begin{tcolorbox}[
    blanker,
    sidebyside,
    sidebyside align=top,
    lefthand width=0.55\textwidth,
    top=4pt, bottom=4pt, left=0pt, right=0pt,
    middle=16pt,
    underlay={
        \draw[dashed, black!80, line width=0.8pt, dash pattern=on 4pt off 3pt]
            ([xshift=0.55\textwidth+8pt]frame.north west) --
            ([xshift=0.55\textwidth+8pt]frame.south west);
    }
]
\textbf{Q.18} \quad {\footnotesize\textsc{[NAT $\cdot$ 1 Mark(s) $\cdot$ No Neg]}} \hfill {\footnotesize\textbf{GATE 2024 CS2} \quad \tiny\texttt{2024\_CS2\_CS\_Q35}}
\vspace{2pt}\hrule height 0.3pt\vspace{4pt}

Let $A$ be an array containing integer values. The distance of $A$ is defined as the minimum number of elements in $A$ that must be replaced with another integer so that the resulting array is sorted in non-decreasing order. The distance of the array $[2, 5, 3, 1, 4, 2, 6]$ is \_\_\_\_\_\_\. (Answer in integer)

\par\vspace{6pt}\noindent\textbf{Your Answer:} \fbox{\rule{0pt}{1.3ex}\hspace{3cm}}

\tcblower

\textbf{Official Key:} \textbf{\large 3 to 3} \hfill \textit{\small [Easy]}
\vspace{2pt}\hrule height 0.3pt\vspace{3pt}
\textbf{Topic:} Algorithms \\
\textbf{Subtopic:} \textit{Searching, sorting, hashing}
\vspace{2pt}\hrule height 0.3pt\vspace{3pt}
\textbf{Guide / Hint:} \\
{\small To minimize replacements to make array $A$ sorted in non-decreasing order, keep the longest non-decreasing subsequence (LNDS). For $A = [2, 5, 3, 1, 4, 2, 6]$, a longest non-decreasing subsequence is $[2, 3, 4, 6]$ of length 4. The minimum elements to replace is $n - \text{length(LNDS)} = 7 - 4 = 3$.}

\end{tcolorbox}
\vspace{6pt}\hrule height 0.35pt\vspace{6pt}
\noindent
\begin{tcolorbox}[
    blanker,
    sidebyside,
    sidebyside align=top,
    lefthand width=0.55\textwidth,
    top=4pt, bottom=4pt, left=0pt, right=0pt,
    middle=16pt,
    underlay={
        \draw[dashed, black!80, line width=0.8pt, dash pattern=on 4pt off 3pt]
            ([xshift=0.55\textwidth+8pt]frame.north west) --
            ([xshift=0.55\textwidth+8pt]frame.south west);
    }
]
\textbf{Q.19} \quad {\footnotesize\textsc{[MCQ $\cdot$ 1 Mark(s) $\cdot$ Neg: 0.33]}} \hfill {\footnotesize\textbf{GATE 2025 CS1} \quad \tiny\texttt{2025\_CS1\_CS\_Q18}}
\vspace{2pt}\hrule height 0.3pt\vspace{4pt}

Let $G$ be any undirected graph with positive edge weights, and $T$ be a minimum spanning tree of $G$. For any two vertices, $u$ and $v$, let $d_1(u, v)$ and $d_2(u, v)$ be the shortest distances between $u$ and $v$ in $G$ and $T$, respectively. Which ONE of the options is CORRECT for all possible $G, T, u$ and $v$?

\par\vspace{6pt}\noindent
\begin{tabular}{@{}p{0.48\linewidth}@{\hspace{0.04\linewidth}}p{0.48\linewidth}@{}}
$\square$\ \textbf{(A)} $d_1(u, v) = d_2(u, v)$ & $\square$\ \textbf{(B)} $d_1(u, v) \le d_2(u, v)$ \\
\addlinespace[3pt]
$\square$\ \textbf{(C)} $d_1(u, v) \ge d_2(u, v)$ & $\square$\ \textbf{(D)} $d_1(u, v) \ne d_2(u, v)$
\end{tabular}

\tcblower

\textbf{Official Key:} \textbf{\large B} \hfill \textit{\small [Easy]}
\vspace{2pt}\hrule height 0.3pt\vspace{3pt}
\textbf{Topic:} Algorithms \\
\textbf{Subtopic:} \textit{Graph traversals, minimum spanning trees, shortest paths}
\vspace{2pt}\hrule height 0.3pt\vspace{3pt}
\textbf{Guide / Hint:} \\
{\small Since the minimum spanning tree $T$ is a subgraph of $G$ ($E(T) \subseteq E(G)$), every path between $u$ and $v$ in $T$ is also a valid path in $G$. Thus, the shortest path distance in $G$ cannot exceed that in $T$, guaranteeing $d_1(u, v) \le d_2(u, v)$.}

\end{tcolorbox}
\vspace{6pt}\hrule height 0.35pt\vspace{6pt}
\noindent
\begin{tcolorbox}[
    blanker,
    sidebyside,
    sidebyside align=top,
    lefthand width=0.55\textwidth,
    top=4pt, bottom=4pt, left=0pt, right=0pt,
    middle=16pt,
    underlay={
        \draw[dashed, black!80, line width=0.8pt, dash pattern=on 4pt off 3pt]
            ([xshift=0.55\textwidth+8pt]frame.north west) --
            ([xshift=0.55\textwidth+8pt]frame.south west);
    }
]
\textbf{Q.20} \quad {\footnotesize\textsc{[MCQ $\cdot$ 1 Mark(s) $\cdot$ Neg: 0.33]}} \hfill {\footnotesize\textbf{GATE 2025 CS1} \quad \tiny\texttt{2025\_CS1\_CS\_Q20}}
\vspace{2pt}\hrule height 0.3pt\vspace{4pt}

Consider the following recurrence relation:
\[ T(n) = 2T(n-1) + n2^n \quad \text{for } n > 0, \; T(0) = 1. \]
Which ONE of the following options is CORRECT?

\par\vspace{6pt}\noindent
\begin{tabular}{@{}p{0.48\linewidth}@{\hspace{0.04\linewidth}}p{0.48\linewidth}@{}}
$\square$\ \textbf{(A)} $T(n) = \Theta(n^2 2^n)$ & $\square$\ \textbf{(B)} $T(n) = \Theta(n 2^n)$ \\
\addlinespace[3pt]
$\square$\ \textbf{(C)} $T(n) = \Theta((\log n)^2 2^n)$ & $\square$\ \textbf{(D)} $T(n) = \Theta(4^n)$
\end{tabular}

\tcblower

\textbf{Official Key:} \textbf{\large A} \hfill \textit{\small [Medium]}
\vspace{2pt}\hrule height 0.3pt\vspace{3pt}
\textbf{Topic:} Algorithms \\
\textbf{Subtopic:} \textit{Asymptotic worst case time and space complexity}
\vspace{2pt}\hrule height 0.3pt\vspace{3pt}
\textbf{Guide / Hint:} \\
{\small Dividing both sides by $2^n$ yields $\frac{T(n)}{2^n} = \frac{T(n-1)}{2^{n-1}} + n$. Letting $S(n) = \frac{T(n)}{2^n}$, we have $S(n) = S(n-1) + n = S(0) + \sum_{i=1}^n i = 1 + \frac{n(n+1)}{2} = \Theta(n^2)$. Multiplying back gives $T(n) = \Theta(n^2 2^n)$.}

\end{tcolorbox}
\vspace{6pt}\hrule height 0.35pt\vspace{6pt}
\noindent
\begin{tcolorbox}[
    blanker,
    sidebyside,
    sidebyside align=top,
    lefthand width=0.55\textwidth,
    top=4pt, bottom=4pt, left=0pt, right=0pt,
    middle=16pt,
    underlay={
        \draw[dashed, black!80, line width=0.8pt, dash pattern=on 4pt off 3pt]
            ([xshift=0.55\textwidth+8pt]frame.north west) --
            ([xshift=0.55\textwidth+8pt]frame.south west);
    }
]
\textbf{Q.21} \quad {\footnotesize\textsc{[NAT $\cdot$ 1 Mark(s) $\cdot$ No Neg]}} \hfill {\footnotesize\textbf{GATE 2025 CS1} \quad \tiny\texttt{2025\_CS1\_CS\_Q33}}
\vspace{2pt}\hrule height 0.3pt\vspace{4pt}

The pseudocode of a function fun() is given below:
\begin{verbatim}
fun(int A[0,...,n-1]){
    for i=0 to n-2
        for j=0 to n-i-2
            if (A[j] > A[j+1])
                then swap A[j] and A[j+1]
}
\end{verbatim}
Let $A[0, \ldots, 29]$ be an array storing 30 distinct integers in descending order. The number of swap operations that will be performed, if the function fun() is called with $A[0, \ldots, 29]$ as argument, is \_\_\_\_\_\_\. (Answer in integer)

\par\vspace{6pt}\noindent\textbf{Your Answer:} \fbox{\rule{0pt}{1.3ex}\hspace{3cm}}

\tcblower

\textbf{Official Key:} \textbf{\large 435 to 435} \hfill \textit{\small [Easy]}
\vspace{2pt}\hrule height 0.3pt\vspace{3pt}
\textbf{Topic:} Algorithms \\
\textbf{Subtopic:} \textit{Searching, sorting, hashing}
\vspace{2pt}\hrule height 0.3pt\vspace{3pt}
\textbf{Guide / Hint:} \\
{\small The function implements standard Bubble Sort. When sorting an array of $n = 30$ elements in strictly reverse order, every comparison triggers a swap operation. The total number of swaps is $\sum_{i=0}^{n-2} (n - 1 - i) = \frac{n(n-1)}{2} = \frac{30 \times 29}{2} = 435$.}

\end{tcolorbox}
\vspace{6pt}\hrule height 0.35pt\vspace{6pt}
\noindent
\begin{tcolorbox}[
    blanker,
    sidebyside,
    sidebyside align=top,
    lefthand width=0.55\textwidth,
    top=4pt, bottom=4pt, left=0pt, right=0pt,
    middle=16pt,
    underlay={
        \draw[dashed, black!80, line width=0.8pt, dash pattern=on 4pt off 3pt]
            ([xshift=0.55\textwidth+8pt]frame.north west) --
            ([xshift=0.55\textwidth+8pt]frame.south west);
    }
]
\textbf{Q.22} \quad {\footnotesize\textsc{[MCQ $\cdot$ 2 Mark(s) $\cdot$ Neg: 0.66]}} \hfill {\footnotesize\textbf{GATE 2025 CS1} \quad \tiny\texttt{2025\_CS1\_CS\_Q43}}
\vspace{2pt}\hrule height 0.3pt\vspace{4pt}

Let $G(V, E)$ be an undirected and unweighted graph with 100 vertices. Let $d(u, v)$ denote the number of edges in a shortest path between vertices $u$ and $v$ in $V$. Let the maximum value of $d(u, v), u, v \in V$ such that $u \neq v$, be 30. Let $T$ be any breadth-first-search tree of $G$. Which ONE of the given options is CORRECT for every such graph $G$?

\par\vspace{6pt}\noindent
\begin{tabular}{@{}p{0.98\linewidth}@{}}
$\square$\ \textbf{(A)} The height of $T$ is exactly 15. \\
\addlinespace[3pt]
$\square$\ \textbf{(B)} The height of $T$ is exactly 30. \\
\addlinespace[3pt]
$\square$\ \textbf{(C)} The height of $T$ is at least 15. \\
\addlinespace[3pt]
$\square$\ \textbf{(D)} The height of $T$ is at least 30.
\end{tabular}

\tcblower

\textbf{Official Key:} \textbf{\large C} \hfill \textit{\small [Medium]}
\vspace{2pt}\hrule height 0.3pt\vspace{3pt}
\textbf{Topic:} Algorithms \\
\textbf{Subtopic:} \textit{Graph traversals, minimum spanning trees, shortest paths}
\vspace{2pt}\hrule height 0.3pt\vspace{3pt}
\textbf{Guide / Hint:} \\
{\small The diameter of $G$ is $\text{diam}(G) = \max_{u,v} d(u, v) = 30$. For any BFS root $r$, the triangle inequality ensures that for the pair $(u, v)$ realizing the diameter, $30 = d(u, v) \le d(r, u) + d(r, v) \le 2 \times \text{height}(T)$. Hence, $\text{height}(T) \ge \lceil 30/2 \rceil = 15$.}

\end{tcolorbox}
\vspace{6pt}\hrule height 0.35pt\vspace{6pt}
\noindent
\begin{tcolorbox}[
    blanker,
    sidebyside,
    sidebyside align=top,
    lefthand width=0.55\textwidth,
    top=4pt, bottom=4pt, left=0pt, right=0pt,
    middle=16pt,
    underlay={
        \draw[dashed, black!80, line width=0.8pt, dash pattern=on 4pt off 3pt]
            ([xshift=0.55\textwidth+8pt]frame.north west) --
            ([xshift=0.55\textwidth+8pt]frame.south west);
    }
]
\textbf{Q.23} \quad {\footnotesize\textsc{[NAT $\cdot$ 2 Mark(s) $\cdot$ No Neg]}} \hfill {\footnotesize\textbf{GATE 2025 CS1} \quad \tiny\texttt{2025\_CS1\_CS\_Q64}}
\vspace{2pt}\hrule height 0.3pt\vspace{4pt}

The maximum value of $x$ such that the edge between the nodes B and C is included in every minimum spanning tree of the given graph is \_\_\_\_\_\_\. (answer in integer)

\par\vspace{6pt}\noindent\textbf{Your Answer:} \fbox{\rule{0pt}{1.3ex}\hspace{3cm}}

\tcblower

\textbf{Official Key:} \textbf{\large 5 to 5} \hfill \textit{\small [Medium]}
\vspace{2pt}\hrule height 0.3pt\vspace{3pt}
\textbf{Topic:} Algorithms \\
\textbf{Subtopic:} \textit{Graph traversals, minimum spanning trees, shortest paths}
\vspace{2pt}\hrule height 0.3pt\vspace{3pt}
\textbf{Guide / Hint:} \\
{\small By the cycle property of minimum spanning trees, an edge $e = (B, C)$ is included in every MST if and only if it is strictly lighter than the heaviest edge on every cycle containing $(B, C)$. Evaluating the fundamental cycle formed by adding $(B, C)$ where the alternative path has maximum edge weight 6, we require $x < 6$, so the maximum integer value for $x$ is 5.}

\end{tcolorbox}
\vspace{6pt}\hrule height 0.35pt\vspace{6pt}
\noindent
\begin{tcolorbox}[
    blanker,
    sidebyside,
    sidebyside align=top,
    lefthand width=0.55\textwidth,
    top=4pt, bottom=4pt, left=0pt, right=0pt,
    middle=16pt,
    underlay={
        \draw[dashed, black!80, line width=0.8pt, dash pattern=on 4pt off 3pt]
            ([xshift=0.55\textwidth+8pt]frame.north west) --
            ([xshift=0.55\textwidth+8pt]frame.south west);
    }
]
\textbf{Q.24} \quad {\footnotesize\textsc{[NAT $\cdot$ 2 Mark(s) $\cdot$ No Neg]}} \hfill {\footnotesize\textbf{GATE 2025 CS1} \quad \tiny\texttt{2025\_CS1\_CS\_Q65}}
\vspace{2pt}\hrule height 0.3pt\vspace{4pt}

In a double hashing scheme, $h_1(k) = k \bmod 11$ and $h_2(k) = 1 + (k \bmod 7)$ are the auxiliary hash functions. The size $m$ of the hash table is 11. The hash function for the $i$-th probe in the open address table is $[h_1(k) + i \cdot h_2(k)] \bmod m$. The following keys are inserted in the given order: 63, 50, 25, 79, 67, 24.
The slot at which key 24 gets stored is \_\_\_\_\_\_\. (Answer in integer)

\par\vspace{6pt}\noindent\textbf{Your Answer:} \fbox{\rule{0pt}{1.3ex}\hspace{3cm}}

\tcblower

\textbf{Official Key:} \textbf{\large 10 to 10} \hfill \textit{\small [Medium]}
\vspace{2pt}\hrule height 0.3pt\vspace{3pt}
\textbf{Topic:} Algorithms \\
\textbf{Subtopic:} \textit{Searching, sorting, hashing}
\vspace{2pt}\hrule height 0.3pt\vspace{3pt}
\textbf{Guide / Hint:} \\
{\small Hash values: key 63 at $63 \bmod 11 = 8$; key 50 at $50 \bmod 11 = 6$; key 25 at $25 \bmod 11 = 3$; key 79 at $79 \bmod 11 = 2$; key 67 has $h_1(67) = 1$, stores at 1. For key 24: $h_1(24) = 24 \bmod 11 = 2$ (collision with 79), $h_2(24) = 1 + (24 \bmod 7) = 1 + 3 = 4$. Probe $i=1$: $(2 + 4) \bmod 11 = 6$ (collision with 50). Probe $i=2$: $(2 + 8) \bmod 11 = 10$ (empty). Hence, key 24 is stored in slot 10.}

\end{tcolorbox}
\vspace{6pt}\hrule height 0.35pt\vspace{6pt}
\noindent
\begin{tcolorbox}[
    blanker,
    sidebyside,
    sidebyside align=top,
    lefthand width=0.55\textwidth,
    top=4pt, bottom=4pt, left=0pt, right=0pt,
    middle=16pt,
    underlay={
        \draw[dashed, black!80, line width=0.8pt, dash pattern=on 4pt off 3pt]
            ([xshift=0.55\textwidth+8pt]frame.north west) --
            ([xshift=0.55\textwidth+8pt]frame.south west);
    }
]
\textbf{Q.25} \quad {\footnotesize\textsc{[MCQ $\cdot$ 1 Mark(s) $\cdot$ Neg: 0.33]}} \hfill {\footnotesize\textbf{GATE 2025 CS2} \quad \tiny\texttt{2025\_CS2\_CS\_Q20}}
\vspace{2pt}\hrule height 0.3pt\vspace{4pt}

Consider an unordered list of $N$ distinct integers. What is the minimum number of element comparisons required to find an integer in the list that is NOT the largest in the list?

\par\vspace{6pt}\noindent
\begin{tabular}{@{}p{0.48\linewidth}@{\hspace{0.04\linewidth}}p{0.48\linewidth}@{}}
$\square$\ \textbf{(A)} 1 & $\square$\ \textbf{(B)} $N - 1$ \\
\addlinespace[3pt]
$\square$\ \textbf{(C)} $N$ & $\square$\ \textbf{(D)} $2N - 1$
\end{tabular}

\tcblower

\textbf{Official Key:} \textbf{\large A} \hfill \textit{\small [Easy]}
\vspace{2pt}\hrule height 0.3pt\vspace{3pt}
\textbf{Topic:} Algorithms \\
\textbf{Subtopic:} \textit{Searching and comparison complexity}
\vspace{2pt}\hrule height 0.3pt\vspace{3pt}
\textbf{Guide / Hint:} \\
{\small To find any element that is not the maximum, compare any two distinct elements in 1 comparison; the smaller of the two is guaranteed not to be the maximum.}

\end{tcolorbox}
\vspace{6pt}\hrule height 0.35pt\vspace{6pt}
\noindent
\begin{tcolorbox}[
    blanker,
    sidebyside,
    sidebyside align=top,
    lefthand width=0.55\textwidth,
    top=4pt, bottom=4pt, left=0pt, right=0pt,
    middle=16pt,
    underlay={
        \draw[dashed, black!80, line width=0.8pt, dash pattern=on 4pt off 3pt]
            ([xshift=0.55\textwidth+8pt]frame.north west) --
            ([xshift=0.55\textwidth+8pt]frame.south west);
    }
]
\textbf{Q.26} \quad {\footnotesize\textsc{[MSQ $\cdot$ 1 Mark(s) $\cdot$ No Neg]}} \hfill {\footnotesize\textbf{GATE 2025 CS2} \quad \tiny\texttt{2025\_CS2\_CS\_Q29}}
\vspace{2pt}\hrule height 0.3pt\vspace{4pt}

Which of the following statements regarding Breadth First Search (BFS) and Depth First Search (DFS) on an undirected simple graph $G$ is/are TRUE?

\par\vspace{6pt}\noindent
\begin{tabular}{@{}p{0.98\linewidth}@{}}
$\square$\ \textbf{(A)} A DFS tree of $G$ is a Shortest Path tree of $G$. \\
\addlinespace[3pt]
$\square$\ \textbf{(B)} Every non-tree edge of $G$ with respect to a DFS tree is a forward/back edge. \\
\addlinespace[3pt]
$\square$\ \textbf{(C)} If $(u, v)$ is a non-tree edge of $G$ with respect to a BFS tree, then the distances from the source vertex $s$ to $u$ and $v$ in the BFS tree are within $\pm 1$ of each other. \\
\addlinespace[3pt]
$\square$\ \textbf{(D)} Both BFS and DFS can be used to find the connected components of $G$.
\end{tabular}

\tcblower

\textbf{Official Key:} \textbf{\large B;C;D} \hfill \textit{\small [Medium]}
\vspace{2pt}\hrule height 0.3pt\vspace{3pt}
\textbf{Topic:} Algorithms \\
\textbf{Subtopic:} \textit{Graph traversals: BFS and DFS}
\vspace{2pt}\hrule height 0.3pt\vspace{3pt}
\textbf{Guide / Hint:} \\
{\small In an undirected graph, non-tree DFS edges can only be back edges (B is true). In a BFS tree, cross edges can only exist between the same or adjacent levels, so distances differ by at most $\pm 1$ (C is true). Both can compute connected components (D is true).}

\end{tcolorbox}
\vspace{6pt}\hrule height 0.35pt\vspace{6pt}
\noindent
\begin{tcolorbox}[
    blanker,
    sidebyside,
    sidebyside align=top,
    lefthand width=0.55\textwidth,
    top=4pt, bottom=4pt, left=0pt, right=0pt,
    middle=16pt,
    underlay={
        \draw[dashed, black!80, line width=0.8pt, dash pattern=on 4pt off 3pt]
            ([xshift=0.55\textwidth+8pt]frame.north west) --
            ([xshift=0.55\textwidth+8pt]frame.south west);
    }
]
\textbf{Q.27} \quad {\footnotesize\textsc{[MCQ $\cdot$ 2 Mark(s) $\cdot$ Neg: 0.66]}} \hfill {\footnotesize\textbf{GATE 2025 CS2} \quad \tiny\texttt{2025\_CS2\_CS\_Q37}}
\vspace{2pt}\hrule height 0.3pt\vspace{4pt}

Let $G$ be an edge-weighted undirected graph with positive edge weights. Suppose a positive constant $\alpha$ is added to the weight of every edge in $G$. Which ONE of the following statements is TRUE?

\par\vspace{6pt}\noindent
\begin{tabular}{@{}p{0.98\linewidth}@{}}
$\square$\ \textbf{(A)} Every MST remains an MST, and every SP remains an SP. \\
\addlinespace[3pt]
$\square$\ \textbf{(B)} MSTs need not remain MSTs, and every SP remains an SP. \\
\addlinespace[3pt]
$\square$\ \textbf{(C)} Every MST remains an MST, and SPs need not remain SPs. \\
\addlinespace[3pt]
$\square$\ \textbf{(D)} MSTs need not remain MSTs, and SPs need not remain SPs.
\end{tabular}

\tcblower

\textbf{Official Key:} \textbf{\large C} \hfill \textit{\small [Medium]}
\vspace{2pt}\hrule height 0.3pt\vspace{3pt}
\textbf{Topic:} Algorithms \\
\textbf{Subtopic:} \textit{Graph algorithms: minimum spanning trees and shortest paths}
\vspace{2pt}\hrule height 0.3pt\vspace{3pt}
\textbf{Guide / Hint:} \\
{\small Adding a constant $\alpha$ to all edge weights preserves the relative order of all spanning trees because every spanning tree has exactly $V - 1$ edges (so all tree costs increase by $(V-1)\alpha$). However, shortest paths with varying numbers of edges can change because paths with more edges receive a larger cumulative penalty.}

\end{tcolorbox}
\vspace{6pt}\hrule height 0.35pt\vspace{6pt}
\noindent
\begin{tcolorbox}[
    blanker,
    sidebyside,
    sidebyside align=top,
    lefthand width=0.55\textwidth,
    top=4pt, bottom=4pt, left=0pt, right=0pt,
    middle=16pt,
    underlay={
        \draw[dashed, black!80, line width=0.8pt, dash pattern=on 4pt off 3pt]
            ([xshift=0.55\textwidth+8pt]frame.north west) --
            ([xshift=0.55\textwidth+8pt]frame.south west);
    }
]
\textbf{Q.28} \quad {\footnotesize\textsc{[MCQ $\cdot$ 2 Mark(s) $\cdot$ Neg: 0.66]}} \hfill {\footnotesize\textbf{GATE 2025 CS2} \quad \tiny\texttt{2025\_CS2\_CS\_Q41}}
\vspace{2pt}\hrule height 0.3pt\vspace{4pt}

An array $A$ of length $n$ with distinct elements is said to be bitonic if there is an index $1 \le i \le n$ such that $A[1..i]$ is strictly increasing and $A[i..n]$ is strictly decreasing. What is the optimal time complexity to search for a key in a bitonic array?

\par\vspace{6pt}\noindent
\begin{tabular}{@{}p{0.48\linewidth}@{\hspace{0.04\linewidth}}p{0.48\linewidth}@{}}
$\square$\ \textbf{(A)} $\Theta(n)$ & $\square$\ \textbf{(B)} $\Theta(1)$ \\
\addlinespace[3pt]
$\square$\ \textbf{(C)} $\Theta(\log^2 n)$ & $\square$\ \textbf{(D)} $\Theta(\log n)$
\end{tabular}

\tcblower

\textbf{Official Key:} \textbf{\large D} \hfill \textit{\small [Medium]}
\vspace{2pt}\hrule height 0.3pt\vspace{3pt}
\textbf{Topic:} Algorithms \\
\textbf{Subtopic:} \textit{Searching and divide-and-conquer: bitonic array search}
\vspace{2pt}\hrule height 0.3pt\vspace{3pt}
\textbf{Guide / Hint:} \\
{\small In a bitonic array, the maximum element can be found in $\Theta(\log n)$ time using binary search on mid-elements. Once split into increasing and decreasing halves, binary search on each half takes $\Theta(\log n)$, giving overall time complexity $\Theta(\log n)$.}

\end{tcolorbox}
\vspace{6pt}\hrule height 0.35pt\vspace{6pt}
\noindent
\begin{tcolorbox}[
    blanker,
    sidebyside,
    sidebyside align=top,
    lefthand width=0.55\textwidth,
    top=4pt, bottom=4pt, left=0pt, right=0pt,
    middle=16pt,
    underlay={
        \draw[dashed, black!80, line width=0.8pt, dash pattern=on 4pt off 3pt]
            ([xshift=0.55\textwidth+8pt]frame.north west) --
            ([xshift=0.55\textwidth+8pt]frame.south west);
    }
]
\textbf{Q.29} \quad {\footnotesize\textsc{[NAT $\cdot$ 2 Mark(s) $\cdot$ No Neg]}} \hfill {\footnotesize\textbf{GATE 2025 CS2} \quad \tiny\texttt{2025\_CS2\_CS\_Q59}}
\vspace{2pt}\hrule height 0.3pt\vspace{4pt}

Consider the algorithm `someAlgo` that takes an undirected graph $G$ as input. What is the return value of `someAlgo(G)`?

\par\vspace{6pt}\noindent\textbf{Your Answer:} \fbox{\rule{0pt}{1.3ex}\hspace{3cm}}

\tcblower

\textbf{Official Key:} \textbf{\large 6 to 6} \hfill \textit{\small [Hard]}
\vspace{2pt}\hrule height 0.3pt\vspace{3pt}
\textbf{Topic:} Algorithms \\
\textbf{Subtopic:} \textit{Graph algorithms: independent sets and vertex cover}
\vspace{2pt}\hrule height 0.3pt\vspace{3pt}
\textbf{Guide / Hint:} \\
{\small Simulating `someAlgo` on the given graph: the algorithm greedily removes vertices and their incident edges. The number of iterations or resulting set cardinality evaluates to 6.}

\end{tcolorbox}
\vspace{6pt}\hrule height 0.35pt\vspace{6pt}
\end{document}